\section{Vérifications et limites}\label{sec:verifications}
\subsection{Vérifications indépendantes}
Les scripts relisent les Parquet avec Pandas sans utiliser les comptes de MongoDB, puis comparent
leur calcul à la réponse de l'API. Les fichiers sont dans \chemin{results/verification/}.
\begin{itemize}
\item \textbf{Janvier, base temporaire.} \code{verify-f2} à \code{verify-f5-parquet.py} réussissent
sur F2 (zone 181, Brooklyn avec 25\,258 candidats et 24\,860 retenus), F3 (zone 132, Brooklyn avec
3 cellules inhabituelles sur 610), F4 (132 $\to$ 161, 161 $\to$ 236, échantillon insuffisant pour
181 $\to$ 89) et F5 (zone 132, Brooklyn, ville avec 19\,368 paires). Nous avons lancé ce contrôle le 7 novembre, lors de l'essai complet sur janvier et avril. Le code de l'API et de l'ingestion est
le même que dans la version finale, seuls des scripts ayant changé entre les deux.
\item \textbf{Six mois, base complète.} \code{verify-registered-period.py} retrouve F1 sur la zone
181 à l'arrivée (50\,064 candidats, 49\,616 retenus), F2 sur Brooklyn de janvier à mars (97\,219 et
95\,809), F4 de JFK à Midtown autour du 10 mars (1\,792 et 2\,220 trajets, médianes 2\,963,5 et
3\,086,5 s) et les trajets d'avril vus par F6.
\item \textbf{Traçabilité.} \code{verify-trips.py} vérifie que \code{/api/trips} renvoie pour la zone
181 les mêmes comptes que F1 (983 candidats, 977 retenus) avec le bon fichier source.
\item \textbf{Changement d'heure.} \code{verify-f4-dst.py} retrouve les 10\,891 s de la ligne
942\,762 de mars et la durée absente de la ligne 841\,170. Il porte sur ces deux lignes réelles
dans une base temporaire et ne valide pas le mois entier.
\end{itemize}
Les archives de chaque benchmark de la section~\ref{sec:mesures} enregistrent aussi des réponses
JSON identiques avant et après index.

\subsection{Limites}
\begin{itemize}
\item Tous les processus partagent le CPU, la mémoire et le disque d'un même PC. Les gains de E2 et
E3 ne se transposent pas à des machines indépendantes, où le réseau compterait et où les disques ne
seraient plus partagés.
\item Les mesures de E1 et E2 reposent sur cinq appels et celles de E3 sur un seul passage. Jenkins
tournait pendant les mesures.
\item Les mesures de la section~\ref{sec:mesures} viennent d'un serveur Linux à deux CPU, sur janvier
et avant les optimisations de F3 et F5. Elles justifient les index, mais ne se comparent pas à E1--E3.
L'index \code{pickup_datetime_zone} n'a pas de mesure séparée.
\item Trois causes restent à établir : le surcoût d'index en réparti, le chargement plus lent en
période et les F4 lents sur shard1 dans E3.
\item Les seuils statistiques restent descriptifs, les calendriers sont incomplets pour les événements
locaux et les flags ne détectent pas toute anomalie. 
\end{itemize}
